
\documentclass{article}
\usepackage{colortbl}
\usepackage{makecell}
\usepackage{multirow}
\usepackage{supertabular}

\begin{document}

\newcounter{utterance}

\centering \large Interaction Transcript for game `textmapworld\_graphreasoning', experiment `large', episode 24 with reka{-}flash{-}3.1{-}t0.0.
\vspace{24pt}

{ \footnotesize  \setcounter{utterance}{1}
\setlength{\tabcolsep}{0pt}
\begin{supertabular}{c@{$\;$}|p{.15\linewidth}@{}p{.15\linewidth}p{.15\linewidth}p{.15\linewidth}p{.15\linewidth}p{.15\linewidth}}
    \# & $\;$A & \multicolumn{4}{c}{Game Master} & $\;\:$B\\
    \hline

    \theutterance \stepcounter{utterance}  
    & & \multicolumn{4}{p{0.6\linewidth}}{
        \cellcolor[rgb]{0.9,0.9,0.9}{
            \makecell[{{p{\linewidth}}}]{
                \texttt{\tiny{[P1$\langle$GM]}}
                \texttt{Please help me with the following task. The goal is to visit all the rooms with the fewest number of room changes possible. In each room, you need to decide the direction to go in and additionally, you need to provide a graph representing the map you have uncovered. Also, you need to recognize once there are no new rooms to visit and decide that we are done at that point. Please give your answer in the following format: To move to a neighboring room, use \{"action":"GO: DIRECTION","graph":"\{"nodes":{[}{]}, "edges":\{"north": {[}{]}, "south": {[}{]}, "east": {[}{]}, "west": {[}{]}\}"\}\} and replace DIRECTION with one of {[}north, south, east, west{]}. To stop the exploration, answer with "DONE" instead. Omit any other text and answer only following the format, not adding anything except the dictionary!} \\
\texttt{Here is an example: } \\
\texttt{You are in the Living Room. Currently available directions: south, west. What is your next instruction?} \\
\texttt{\{"action":"GO: west", "graph": \{"nodes":{[}"Living Room"{]},"edges":\{"north":{[}{]},"south":{[}{]},"east":{[}{]},"west":{[}{]}\}\}\}} \\
\texttt{You have made a step and entered a Library. Currently available directions: east, north. What is your next instruction?} \\
\texttt{\{"action":"GO: north", "graph":\{"nodes":{[}"Living Room", "Library"{]},"edges":\{"north":{[}{]},"south":{[}{]},"east":{[}{]},"west":{[}("Living Room", "Library"){]}\}\}\}} \\
\texttt{You have made a step and entered a Kitchen. Currently available directions: south, east. What is your next instruction?} \\
\texttt{\{"action":"GO: east", "graph":\{"nodes": {[}"Living Room", "Library", "Kitchen"{]}, "edges":\{"north": {[}("Library", "Kitchen"){]},"south": {[}{]},"east": {[}{]},"west": {[}("Living Room", "Library"){]}\}\}\}} \\
\texttt{...} \\
\texttt{You have made a step and entered a Bedroom. Currently available directions: south, west. What is your next instruction?} \\
\texttt{\{"action": "DONE", "graph": \{...\}\}} \\
\texttt{Let us start.} \\
\texttt{You are in the Bar. Currently available directions: west. What is your next instruction?} \\
            }
        }
    }
    & \\ \\

    \theutterance \stepcounter{utterance}  
    & \multicolumn{4}{p{0.6\linewidth}}{
        \cellcolor[rgb]{0.9,0.9,0.9}{
            \makecell[{{p{\linewidth}}}]{
                \texttt{\tiny{[P1$\rangle$GM]}}
                \texttt{<reasoning>} \\
\texttt{Let's tackle this problem. The goal is to visit all rooms with the fewest room changes possible. So, I need to figure out the shortest path to explore every room. This sounds like a classic graph traversal problem, probably using BFS (Breadth{-}First Search) since BFS is known for finding the shortest path in unweighted graphs.} \\
\\ 
\texttt{First, the user starts in the Bar, and the only available direction is west. So the first move has to be west. But wait, the example shows that when moving, the graph is updated with edges. Let me look at the example again. In the first step, the user is in the Living Room, goes west to the Library. The graph's edges for west have ("Living Room", "Library"). So edges are bidirectional? Because when moving west from Living Room to Library, the edge is added in the west direction of Living Room and maybe east in Library? Wait, the example's second step shows that after moving north from Library to Kitchen, the edges for north in Library have ("Library", "Kitchen"). But in the graph structure, each node's edges are listed per direction. So maybe each edge is represented once in each direction. For example, moving west from A to B would add B to A's west and A to B's east. But in the example, when moving west from Living Room to Library, the west edge of Living Room is empty, but the west edge of Library isn't mentioned. Wait, looking at the example's first response: the graph's edges for west is empty for Living Room. Wait, the first step's graph is \{"nodes":{[}"Living Room"{]},"edges":\{"north":{[}{]},"south":{[}{]},"east":{[}{]},"west":{[}{]}\}\}. Then after moving west to Library, the graph's edges for west of Living Room is {[}("Living Room", "Library"){]}? Wait, no. Let me check the example again.} \\
\\ 
\texttt{In the first example step:} \\
\\ 
\texttt{User is in Living Room, available directions: south, west. The action is GO: west. The graph is \{"nodes":{[}"Living Room"{]},"edges":\{"north":{[}{]},"south":{[}{]},"east":{[}{]},"west":{[}{]}\}\}. Wait, but after moving west, the next step is in the Library. The graph in the second step is \{"nodes":{[}"Living Room", "Library"{]},"edges":\{"north":{[}{]},"south":{[}{]},"east":{[}{]},"west":{[}("Living Room", "Library"){]}\}\}. So the west edge of} \\
            }
        }
    }
    & & \\ \\

    \theutterance \stepcounter{utterance}  
    & & & \multicolumn{2}{p{0.3\linewidth}}{
        \cellcolor[rgb]{0.9,0.9,0.9}{
            \makecell[{{p{\linewidth}}}]{
                \texttt{\tiny{[GM$|$GM]}}
                \texttt{abort game} \\
            }
        }
    }
    & & \\ \\

\end{supertabular}
}

\end{document}
